\subsection{Algorithmic Comparison}

\Cref{tab:full_performance} presents the primary N=32 comparison. $U_{PRR}$ (\Cref{eq:uprr}) is the packet-delivery training signal; $U_{OTR}$ (\Cref{eq:uotr}) is the deadline-compliance criterion. At dense fleet scale, MA-ROPE improves over AllOffload on $U_{PRR}$ (0.8323 vs. 0.7790), $U_{OTR}$ (0.7424 vs. 0.5644), fleet OTR (0.8340 vs. 0.5250), and energy (0.7057 J vs. 0.9441 J).

\begin{table}[htbp]
\caption{Full Performance Comparison---URLLC stress, N=32 vehicles, 5 seeds. $U_{PRR}$: Eq. (4a), training signal. $U_{OTR}$: Eq. (4b), deadline compliance. Fleet OTR = on-time offload PDUs / transmitted offload PDUs.}
\label{tab:full_performance}
\centering
\resizebox{\columnwidth}{!}{%
\begin{tabular}{lrrrrr}
\toprule
\textbf{Method} & \textbf{U\_PRR} & \textbf{U\_OTR} & \textbf{Fleet OTR} & \textbf{Energy (J)} & \textbf{Latency (ms)} \\
\midrule
\textbf{MA-ROPE Full $\star$} & \textbf{0.8323} & \textbf{0.7424} & \textbf{0.8340} & \textbf{0.7057} & \textbf{213.0} \\
Yun2022 (GA+DQN) & 0.8134 & 0.7415 & 0.8340 & 0.7448 & 208.8 \\
Wu2021 (UCB Bandit) & 0.8268 & 0.7437 & 0.8340 & 0.7221 & 204.0 \\
Jing2019 (Lyapunov) & 0.8339 & 0.7414 & 0.8340 & 0.5032 & 275.6 \\
Shao2025 (CoopEC) & 0.8339 & 0.7414 & 0.8340 & 0.5032 & 275.6 \\
DQN --- Shao constraints & 0.8378 & 0.7465 & 0.8340 & 0.6822 & 202.0 \\
DQN --- Yun constraints & 0.8361 & 0.7481 & 0.8340 & 0.7492 & 182.3 \\
DQN --- Wu constraints & 0.8219 & 0.7449 & 0.8340 & 0.8073 & 182.5 \\
DQN --- Jing constraints & 0.8328 & 0.7488 & 0.8340 & 0.7643 & 176.9 \\
MADDPG (upper bound) & 0.8298 & 0.7518 & 0.8340 & 0.9460 & 139.1 \\
AllOffload (5G 100\%) & 0.7790 & 0.5644 & 0.5250 & 0.9441 & 176.9 \\
NoOffload (local) & 0.8442 & 0.2754 & 0.0000 & 0.5032 & 248.1 \\
\bottomrule
\end{tabular}%
}
\end{table}

\begin{figure}[htbp]
\centering
\includegraphics[width=\columnwidth]{example-image-a}
\caption{Utility vs. Vehicle Speed--URLLCStress Scenario.}
\label{fig:utility_speed}
\end{figure}

\begin{figure}[htbp]
\centering
\includegraphics[width=\columnwidth]{example-image-b}
\caption{Average Latency vs. Vehicle Speed--URLLCStress Scenario.}
\label{fig:latency_speed}
\end{figure}

\begin{figure}[htbp]
\centering
\includegraphics[width=\columnwidth]{example-image-c}
\caption{Energy Consumption vs. Vehicle Speed--URLLCStress Scenario.}
\label{fig:energy_speed}
\end{figure}

\begin{figure}[htbp]
\centering
\includegraphics[width=\columnwidth]{example-image-a}
\caption{Empirical PRR = $n_{rx}/n_{tx}$ vs. Vehicle Speed. AllOffload degrades fastest (5G-NR Jakes). MA-ROPE remains stable via ITS-G5 adaptation.}
\label{fig:prr_speed}
\end{figure}

\subsection{Table IV Panels: Literature, Baselines, Restricted DQN, MADDPG, and Safety}

\Cref{tab:lit_replicas,tab:fixed_baselines,tab:restricted_dqn,tab:maddpg_upper,tab:safety_task} decompose \Cref{tab:full_performance} into the reviewer-facing comparisons: literature replicas, fixed baselines, restricted DQN action spaces, the MADDPG upper bound, and safety-only OTR at N=32.

\begin{table}[htbp]
\caption{Literature Replicas vs. MA-ROPE--URLLC stress, N=32 vehicles}
\label{tab:lit_replicas}
\centering
\resizebox{\columnwidth}{!}{%
\begin{tabular}{lrrrrr}
\toprule
\textbf{Method} & \textbf{U\_PRR} & \textbf{U\_OTR} & \textbf{Fleet OTR} & \textbf{Energy (J)} & \textbf{Latency (ms)} \\
\midrule
\textbf{MA-ROPE Full $\star$} & \textbf{0.8323} & \textbf{0.7424} & \textbf{0.8340} & \textbf{0.7057} & \textbf{213.0} \\
Yun2022 (GA+DQN) & 0.8134 & 0.7415 & 0.8340 & 0.7448 & 208.8 \\
Wu2021 (UCB Bandit) & 0.8268 & 0.7437 & 0.8340 & 0.7221 & 204.0 \\
Jing2019 (Lyapunov) & 0.8339 & 0.7414 & 0.8340 & 0.5032 & 275.6 \\
Shao2025 (CoopEC) & 0.8339 & 0.7414 & 0.8340 & 0.5032 & 275.6 \\
\bottomrule
\end{tabular}%
}
\end{table}

\begin{figure}[htbp]
\centering
\includegraphics[width=\columnwidth]{example-image-b}
\caption{Restricted-action DQN comparisons. The figure complements \Cref{tab:restricted_dqn} by showing how action-space constraints affect the learned policy.}
\label{fig:restricted_dqn_plot}
\end{figure}

\begin{table}[htbp]
\caption{Fixed Baselines vs. MA-ROPE--URLLC stress, N=32 vehicles}
\label{tab:fixed_baselines}
\centering
\resizebox{\columnwidth}{!}{%
\begin{tabular}{lrrrrr}
\toprule
\textbf{Method} & \textbf{U\_PRR} & \textbf{U\_OTR} & \textbf{Fleet OTR} & \textbf{Energy (J)} & \textbf{Latency (ms)} \\
\midrule
\textbf{MA-ROPE Full $\star$} & \textbf{0.8323} & \textbf{0.7424} & \textbf{0.8340} & \textbf{0.7057} & \textbf{213.0} \\
AllOffload (5G 100\%) & 0.7790 & 0.5644 & 0.5250 & 0.9441 & 176.9 \\
NoOffload (local) & 0.8442 & 0.2754 & 0.0000 & 0.5032 & 248.1 \\
\bottomrule
\end{tabular}%
}
\end{table}

\begin{table}[htbp]
\caption{Restricted DQN Action Spaces vs. MA-ROPE--URLLC stress, N=32 vehicles}
\label{tab:restricted_dqn}
\centering
\resizebox{\columnwidth}{!}{%
\begin{tabular}{lrrrrr}
\toprule
\textbf{Method} & \textbf{U\_PRR} & \textbf{U\_OTR} & \textbf{Fleet OTR} & \textbf{Energy (J)} & \textbf{Latency (ms)} \\
\midrule
\textbf{MA-ROPE Full $\star$} & \textbf{0.8323} & \textbf{0.7424} & \textbf{0.8340} & \textbf{0.7057} & \textbf{213.0} \\
DQN --- Shao constraints & 0.8378 & 0.7465 & 0.8340 & 0.6822 & 202.0 \\
DQN --- Yun constraints & 0.8361 & 0.7481 & 0.8340 & 0.7492 & 182.3 \\
DQN --- Wu constraints & 0.8219 & 0.7449 & 0.8340 & 0.8073 & 182.5 \\
DQN --- Jing constraints & 0.8328 & 0.7488 & 0.8340 & 0.7643 & 176.9 \\
\bottomrule
\end{tabular}%
}
\end{table}

\begin{table}[htbp]
\caption{MADDPG Upper Bound vs. MA-ROPE--URLLC stress, N=32 vehicles}
\label{tab:maddpg_upper}
\centering
\resizebox{\columnwidth}{!}{%
\begin{tabular}{lrrrrr}
\toprule
\textbf{Method} & \textbf{U\_PRR} & \textbf{U\_OTR} & \textbf{Fleet OTR} & \textbf{Energy (J)} & \textbf{Latency (ms)} \\
\midrule
\textbf{MA-ROPE Full $\star$} & \textbf{0.8323} & \textbf{0.7424} & \textbf{0.8340} & \textbf{0.7057} & \textbf{213.0} \\
MADDPG (upper bound) & 0.8298 & 0.7518 & 0.8340 & 0.9460 & 139.1 \\
\bottomrule
\end{tabular}%
}
\end{table}

\begin{table}[htbp]
\caption{Safety-Task Delivery at N=32--Safety OTR and PRR-OTR Gap}
\label{tab:safety_task}
\centering
\begin{tabular}{lrrr}
\toprule
\textbf{Method} & \textbf{Safety PRR} & \textbf{Safety OTR} & \textbf{PRR-OTR gap} \\
\midrule
\textbf{MA-ROPE Full $\star$} & \textbf{0.9890} & \textbf{0.7519} & \textbf{23.7\%} \\
AllOffload & 0.9790 & 0.2869 & 69.2\% \\
Wu2021 & 0.9890 & 0.7519 & 23.7\% \\
Yun2022 & 0.9890 & 0.7520 & 23.7\% \\
MADDPG & 0.9890 & 0.7520 & 23.7\% \\
\bottomrule
\end{tabular}
\end{table}

\subsection{Fleet-Density Scaling and RSU Queue Contention}

\Cref{tab:why_marope_wins} summarizes the deployment story across fleet OTR, $U_{OTR}$, energy, queue loss, and step overhead. The fleet sweep shows the central scaling effect: AllOffload drops from 87.7\% OTR at N=8 to 52.5\% at N=32, while MA-ROPE remains at 83.4\%.

\begin{table}[htbp]
\caption{Why MA-ROPE Wins at Fleet Scale--Utility, Energy, Queue Loss, and OTR}
\label{tab:why_marope_wins}
\centering
\begin{tabular}{lrrl}
\toprule
\textbf{Metric} & \textbf{MA-ROPE} & \textbf{Comparator} & \textbf{Story} \\
\midrule
Fleet OTR @ N=8 & 0.887 & AllOffload 0.877 & +1.0 pp \\
Fleet OTR @ N=16 & 0.869 & AllOffload 0.757 & +11.2 pp \\
Fleet OTR @ N=32 & 0.834 & AllOffload 0.525 & +30.9 pp \\
$U_{OTR}$ @ N=32 & 0.742 & AllOffload 0.564 & +0.178 \\
Energy @ N=32 & 0.706 J & AllOffload 0.944 J & 25.2\% lower \\
Queue loss @ N=32 & 1.15\% & AllOffload 2.14\% & 1.9x lower \\
Step time & 166.6 \textmu s & Single-RAT 167.6 \textmu s & No penalty \\
\bottomrule
\end{tabular}
\end{table}

\begin{figure}[htbp]
\centering
\includegraphics[width=\columnwidth]{example-image-c}
\caption{RSU queue-classified loss fraction vs. fleet size (URLLC, 5 seeds). At N=32, MA-ROPE reports 1.15\% queue loss versus 2.14\% for AllOffload; queue loss is a supporting metric, while fleet OTR carries the stronger scaling story.}
\label{fig:queue_loss}
\end{figure}

\begin{figure}[htbp]
\centering
\includegraphics[width=\columnwidth]{example-image-a}
\caption{System utility vs. fleet size (URLLC, 5 seeds). The MA-ROPE utility gap over AllOffload widens from +0.0277 at N=8 to +0.0533 at N=32.}
\label{fig:system_utility}
\end{figure}

\Cref{fig:fleet_otr} extends the fleet-scaling analysis to on-time delivery. At N=8, MA-ROPE achieves fleet OTR = 0.887 vs. AllOffload = 0.877; at N=16 the gap widens to 0.869 vs. 0.757; and at N=32 it becomes 0.834 vs. 0.525. These values are taken from table\_vehicle\_density.csv and provide the main dense-fleet scaling headline.

\subsection{Discussion}

NoOffload vs. MA-ROPE on composite utility. \Cref{tab:full_performance} shows NoOffload achieving higher $U_{PRR}$ (0.8442) than MA-ROPE (0.8323) at N=32, but NoOffload has no transmitted offload PDUs and fleet OTR is therefore not a meaningful V2X delivery metric. Its $U_{OTR}$ is only 0.2754. The operational claim should therefore be based on fleet OTR, $U_{OTR}$, energy, and dense-fleet scaling rather than scalar $U_{PRR}$ alone.

Dual utility functions. MA-ROPE exposes two utility criteria: $U_{PRR}$ (\Cref{eq:uprr}) for dense packet-delivery reward shaping and $U_{OTR}$ (\Cref{eq:uotr}) for deadline-compliance-dominant deployment assessment. In the primary URLLC evaluation, $U_{OTR}$ is the safety-facing criterion because a received-but-late safety PDU has zero operational value under SAE J3016. $U_{PRR}$ is reported alongside for balanced and compute-heavy scenarios where throughput dominates over strict timing. The dual-criterion design allows MA-ROPE to be evaluated across V2X use cases without changing the underlying simulator interface. Task-adaptive alpha/beta weights are principled rather than tuned: safety alpha=0.70 matches the 65\% safety PDU proportion; compute alpha=0.20 reflects 2000 ms deadline headroom where throughput dominates. A sensitivity analysis of these weights is left for future work.

Limitations and future work. \Cref{tab:deadline_analysis} remains an N=8 MA-ROPE-only deep safety evaluation, while \Cref{tab:safety_task} is the multi-method N=32 safety comparison. Future work should evaluate non-IID vehicle/task distributions, secure or differentially private FedAvg, trace-driven mobility validation, RSU handover/failure scenarios, and sensitivity to alpha/beta weights and queue service rates.

\begin{figure}[htbp]
\centering
\includegraphics[width=\columnwidth]{example-image-b}
\caption{Fleet OTR vs. fleet size (URLLC). MA-ROPE holds 0.834 fleet OTR at N=32, while AllOffload falls to 0.525.}
\label{fig:fleet_otr}
\end{figure}

The results reveal a sharp divergence in deadline compliance. At N=8, MA-ROPE and AllOffload are close in fleet OTR (0.887 vs. 0.877). By N=16, AllOffload falls to 0.757 while MA-ROPE remains at 0.869. At N=32, AllOffload collapses to 0.525, while MA-ROPE maintains 0.834. Queue loss also favors MA-ROPE (1.15\% vs. 2.14\% at N=32), but the decisive deployment story is deadline-compliant fleet delivery at scale.
